% s-208-p208-m002; original PDF 208, 209
\Aphorism{s-208-p208-m002}{300}
Mitleid und Furcht ist für meine Empfindung gerade das, was den \Sperr{untragischen} Situationen im Gegensatz zu den tragischen zukommt. Mit ihnen antworte ich etwa, wenn ich von grausam gemarterten Menschen und Thieren erfahre, aber durchaus nicht, wenn ich Hamlets oder Tristans Schicksal miterlebe. Furcht und Mitleid, selbst in der sublimsten Mischung und Vergeistigung, sind typisch unterschieden von der tragischen Stimmung; entweder fühle ich hierin anders als die moderne Welt, oder wir Modernen anders als die Griechen, oder die Griechen anders als Aristoteles. Wenn Schopenhauer der Tragödie den Zweck unterlegt, den Willen vom Leben\Sperr{ abzuwenden,} so steht er hier ganz auf aristotelischem Grunde und empfindet hedonistisch wie dieser grösste seiner Antipoden; die Isolation des Betrachters vom Kunstwerk, die ästhetische Hülle um den Zuschauer, die für seine persönlichen Alltagsinteressen undurchlässig sein soll, wäre sofort durchbrochen, wenn das Schicksal der Bühnengestalten von Jenem als sein eigenes gefürchtet oder das seiner Nächsten bemitleidet oder als das der menschlichen Gattung verallgemeinert und zum Einwand gegen das Dasein ausgedeutet würde.
\par

